\section{Experimental Setup}
\label{sec:setup}

\subsection{Dataset and split}

AutoViVQA provides 19{,}411 images / 37{,}077 Vietnamese questions, each with five
diverse free-form answers and a reasoning-type label, released with an 80/20
train/val division and no public test split. We re-partition into a
\textbf{grouped 70/15/15 split}: every image is assigned to exactly one of
train / val / test by hashed \texttt{image\_id} (seed 42), so no image --- and
therefore no shared caption or scene context --- crosses splits, and the
reasoning-type distribution is held approximately fixed across splits.

\begin{table}[t]
  \centering
  \caption{Split statistics. Image overlap between any two splits is $0$.}
  \label{tab:splits}
  \begin{tabular}{lrrrr}
    \toprule
    Split & Questions (eval) & Images & Ans/q & Ans length \\
    \midrule
    Train & ${\sim}25{,}900$ & 13{,}576 & 5.00 & 4.3\,w \\
    Val   & 5{,}463 & 2{,}908 & 5.00 & 4.3\,w \\
    Test  & 5{,}468 & 2{,}914 & 5.00 & 4.3\,w \\
    \bottomrule
  \end{tabular}
\end{table}

A random question-level split gives bridge numbers $\le 0.5$ F1 / $\le 3$ CIDEr-D
higher than this grouped split, so prior AutoViVQA bridge results are not badly
inflated by leakage; we adopt the grouped split for all numbers here.

Reasoning-type (\texttt{category}) distribution: \texttt{relational}
(${\sim}30\%$), \texttt{recognition} (${\sim}19\%$), \texttt{spatial}
(${\sim}15\%$), \texttt{causal} (${\sim}13\%$), \texttt{counting} (${\sim}12\%$),
with \texttt{action}, \texttt{context}, \texttt{yes/no} in the tail. We use the
nominal \texttt{category} labels and treat them as unordered.

The oracle sweep (Sec.~\ref{sec:abl-routing}) evaluates every action on every
question, so it uses an equal-per-category subset
($5{,}547$ / $3{,}727$ / $3{,}739$ questions for train / val / test); overall
means are re-weighted to the natural distribution.

\subsection{Answer selection during training}

Each question has five references. Bridge and LoRA training optimise cross-entropy
against the \textbf{first} reference --- a fixed choice, for reproducibility. All
\textbf{evaluation} uses all five references: per-sample max for token metrics,
corpus multi-reference for CIDEr-D / BLEU / ROUGE / METEOR. Resampling a reference
each epoch (\texttt{answer\_sampling = random}) is one of the RQ5 interventions
(Sec.~\ref{sec:abl-rq5}) and is not otherwise used.

\subsection{Training}

InternViT-300M and Qwen2-0.5B are \textbf{frozen}; only the bridge (and, for RQ6,
a rank-16 decoder LoRA) trains. AdamW, learning rate \num{2e-4}, micro-batch 8, no
early stopping. \textbf{Two epochs} for the bridge: validation cross-entropy
bottoms out within one epoch and corpus CIDEr plateaus by epoch 2 (multi-token
CIDEr $1.025$ at epoch 2 $\rightarrow$ $1.029$ at epoch 4 on an 800-sample
probe), so a third and fourth epoch add nothing and risk instability. The decoder
LoRA is trained for 1 or 3 epochs on top of a fixed bridge
(Sec.~\ref{sec:abl-rq6} reports both). One bridge run is ${\sim}5$ GPU-h; one
LoRA run is ${\sim}3$ (1 epoch) to ${\sim}8$ (3 epochs) GPU-h.

\paragraph{Note on a superseded run.}
One early bridge run trained for four epochs; its seed-42 Residual checkpoint was
a training-instability outlier (validation CE $2.35$ vs.\ $1.5$--$1.7$ for every
other bridge and seed). We re-ran all bridges at two epochs and three seeds; the
sound Residual numbers (F1 $45.6$, CIDEr-D $81.1$) are in line with the other
bridges, and all reported numbers use the two-epoch, three-seed runs. An earlier
auxiliary distillation term (MSE between the bridge's first token and a
linear-projector baseline) is disabled everywhere --- it applied to only two of
the five bridges and is an unfair confound.

\subsection{Hardware, seeds, metrics, statistics}
\label{sec:setup-hw}

All training and evaluation ran on Kaggle P100 (16\,GB) and T4 GPU kernels, no
multi-GPU. Total project compute is on the order of a few hundred GPU-hours; a
single headline recipe run (bridge $+$ 3-epoch LoRA) is ${\sim}13$ GPU-h. Every
bridge and negative-axis intervention and the decoder-LoRA rows are
\textbf{three seeds} ($42$, $123$, $3407$; the multi-token headline row adds
$2026$); oracle / routing runs are seed $42$. We report mean $\pm$ standard
deviation.

Cross-paper generation claims use \textbf{corpus CIDEr-D / BLEU-4 / ROUGE-L}
(pycocoevalcap); METEOR is reported but not used for cross-paper claims (it
varies by ${\sim}13$ points between implementations). Answer matching is
token-level P/R/F1 and exact-match, per-sample max over the five references, from
an in-house scorer whose values match the AutoViVQA baseline table
(Table~\ref{tab:main-inhouse} reports both conventions side by side). Main
comparisons use paired bootstrap ($2{,}000$ resamples); we report $95\%$ CIs on
$\Delta$ quantities and $P(\Delta > 0)$. The full validation split is scored
every run; the \textbf{test split is scored once} and never used for model
selection. Libraries: PyTorch~2, Transformers, \texttt{pyvi}, PhoBERT-base-v2;
InternViT-300M / Qwen2-0.5B from Vintern-1B-v3.5.
